\chapter{专长与记忆}\label{s:memory}

\begin{quote}

  记忆是思考的残留物。\\
  —— Daniel Willingham\index{Willingham, Daniel}，\emph{Why Students Don't Like School}

\end{quote}

上一章讲了新手和熟手之间的差别。
这一章来看专长：
它是什么，
人是怎么获得它的，
以及它为何既能帮上忙、又能帮倒忙。
之后我们会介绍学习最重要的一个限制，
并看看画出心智模型的图，如何帮助我们把握知识转化成课程。

先说，我们说某人是专家时，指的是什么？\index{专家}
通常的回答是：他们解决问题比「仅仅胜任」的人快得多，
或者他们能认出并处理那些常规规则不适用的情况。
他们还不知怎的把这一切做得不费吹灰之力：
在许多情况下，
他们似乎一眼就能看出正确答案~\cite{Parn2017}。

专长不只是知道更多事实：
熟手也能记住大量琐碎知识，而表现看不出有什么提升。
不如换个角度，暂且想象我们把知识存成一张网络或图，
其中事实是节点，
关系是弧\footnote{这\emph{绝对}不是我们大脑的工作方式，但它是个有用的比喻。}。
专家和熟手的关键区别在于，
专家的心智模型连接要密集得多，
也就是他们更有可能知道任意两个事实之间的关联。

图这个比喻解释了为什么帮学习者建立连接，
和把事实介绍给他们一样重要：
没有这些连接，
人就无法回忆并运用自己知道的东西。
它也解释了专家行为的许多被观察到的特点：

\begin{itemize}

\item
  专家往往能从一个问题直接跳到答案，
  因为在他们的脑子里，两者之间确实有一条直接连线。
  熟手可能要推演
  $A{\rightarrow}B{\rightarrow}C{\rightarrow}D{\rightarrow}E$，
  专家却能一步从 $A$ 到 $E$。
  我们把这叫做\gref{g:intuition}{直觉}：
  专家不是一路推理出答案，
  而是像认出一张熟脸那样认出答案。

\item
  连接密集的图也是专家\grefdex{g:fluid-representation}{流畅表征}{流畅表征}的基础，
  也就是他们能够在问题的不同视角之间来回切换的能力~\cite{Petr2016}。
  例如，
  在解一道数学题时，
  专家可能一会儿用几何方式处理它，
  一会儿又把它表示成一组方程。

\item
  这个比喻还解释了为什么专家的诊断能力比熟手强：
  事实之间的连线越多，就越容易从症状反推到原因。
  （这也就解释了为什么在求职面试中让程序员调试，
  比让他编程更能准确反映他的能力。）

\item
  最后，
  专家往往对自己的学科熟悉到了
  再也无法想象\emph{不}用那种方式看世界是什么样子。
  这意味着他们教这门学科的能力，往往还不如专业程度较低、
  但仍记得自己当初是怎么学的人。

\end{itemize}

最后这一点叫做\gref{g:expert-blind-spot}{专家盲区}。
按照它在~\cite{Nath2003}里的原始定义，
它指专家倾向于按学科的深层原理来组织讲解，
而不是依据学习者已经知道什么来引导。
它可以通过训练克服，
但它正是这样一个现象的原因之一：
某人在某一领域做研究有多好，
和教这门学科有多好，两者之间没有相关性~\cite{Mars2002}。

\begin{aside}{那个「就」字}
  专家常常会用「就」这个字暴露自己的盲区，
  比如：
  「哦，这很简单，你只要开一台新虚拟机，
  然后只要给 Ubuntu 装这四个补丁，
  然后只要把整个程序用纯函数式语言重写一遍就行了。」
  正如我们在\chapref{s:motivation}里讨论的，
  这么说等于在暗示讲的人觉得这个问题微不足道，
  那么在这个问题上卡住的人想必就是个笨蛋，
  所以别这么干。
\end{aside}

\seclbl{概念图}{s:memory-concept-maps}

我们用来表示某人心智模型的首选工具是\gref{g:concept-map}{概念图}，
其中事实是气泡，连接是带标注的连线。
作为例子，
\figref{f:memory-seasons}展示了地球为什么有四季（来自\hreffoot{https://cmap.ihmc.us/}{IHMC}），
而\appref{s:conceptmaps}从三个视角给出了图书馆的概念图。

\figpdf{figures/seasons.pdf}{四季的概念图}{f:memory-seasons}{}

为了展示概念图在编程教学中怎么用，
看看 Python 里的这个 \texttt{for} 循环：\index{Python}

\begin{minted}{text}
for letter in "abc":
    print(letter)
\end{minted}

\noindent
它的输出是：

\begin{minted}{text}
a
b
c
\end{minted}

这个循环里的三个关键「东西」显示在\figref{f:memory-loop}的上半部分，
但它们只是故事的一半。
下半部分的扩展版展示了这些东西之间的关系，
这些关系和概念本身一样，对理解至关重要。

\figpdf{figures/for-loop.pdf}{一个 \texttt{for} 循环的概念图}{f:memory-loop}{}

\newpage
概念图有很多种用法：

\begin{description}

\item[帮老师弄清自己想教什么。]
  概念图把内容和顺序分开了：
  依我们的经验，
  人们很少会按自己最初画图时的顺序去教。

\item[协助课程设计者之间的沟通。]
  对教什么想法很不一致的老师，
  很可能会把各自的学习者往不同方向拽。
  画概念图并相互分享有助于避免这种情况。
  是的，
  不同的人对同一个主题可能有不同的概念图，
  但概念图的绘制会让这些差异变得明确。

\item[协助与学习者沟通。]
  虽然可以在课程开始时给学习者一张预先画好的图让他们自己标注，
  但更好的做法是一边教一边一块一块地画出来，
  以强化图里的内容与老师说过的内容之间的联系。
  我们会在\secref{s:architecture-brain}回到这个想法。

\item[用于评价。]
  让学习者画出他们觉得自己刚学到的东西，
  能向老师表明他们漏掉了什么、哪里被讲错了。
  评阅学习者的概念图太费时间，没法当作课上的形成性评价，
  但\emph{一旦学习者熟悉了这种技巧}，它在每周的讲授课里就非常有用。
  这个限定是必要的，因为
  任何新的做事方式一开始都会拖慢人——如果学习者正忙着弄懂基础编程，
  同时又让他去琢磨怎么把脑子里想的画出来，那是不公平的负担。

\end{description}

有些老师也怀疑新手能不能有效地把他们的理解画成图，
因为自我反省和解释自己的理解，通常比理解本身是更高级的技能。
例如，
\cite{Kepp2008}考察了概念图在计算机教育中的使用。
他们的发现之一是：
「{\ldots}概念图对许多学生来说很麻烦，因为
它考的是个人的理解，而不是仅仅死记硬背来的知识。」
作为一个看重理解胜过死记硬背的人，
我认为这是好事。

\begin{aside}{从哪儿开始都行}
  第一次被要求画概念图时，很多人不知道该从哪儿下手。
  碰到这种情况，
  写下两个与你打算画的主题相关的词，
  然后在它们之间画一条线，并加一个标注说明这两个想法是怎么关联的。
  接着你可以问：还有哪些东西以同样的方式相关联？
  这些东西各有什么组成部分？
  或者，在纸上已有的概念之前或之后会发生什么？
  这样就能发现更多的节点和弧。
  再往后，难的地方往往是怎么停下来。
\end{aside}

概念图只是我们表示对某一学科理解的一种方式~\cite{Eppl2006}；
其他的还有韦恩图、流程图和决策树~\cite{Abel2009}。
所有这些都\grefdex{g:externalized-cognition}{把认知外化}{把认知外化}，
也就是让心智模型变得可见，从而可以比较和组合\footnote{套用
奥斯卡·王尔德的温德米尔夫人说的话，\index{Wilde, Oscar}
人们往往要听到自己说出口，才知道自己在想什么。}。

\begin{aside}{草稿与诚实}
  许多用户界面设计师认为，
  与其给人看精心打磨的样稿，不如给他们看粗糙的想法草图，
  因为对于自己觉得只花了几分钟做出来的东西，人们更可能给出诚实的反馈：
  如果看起来被批评的东西是花了好几个小时做出来的，
  大多数人就会手下留情。
  所以，在画概念图来激发讨论时，
  你应该用铅笔和废纸（或者笔和白板），
  而不是花哨的计算机绘图工具。
\end{aside}

\seclbl{七加减二}{s:memory-seven-plus-or-minus}

知识的图模型虽然错、但有用，
另一个简单模型却有更可靠的生理基础。
大致而言，
人的记忆可以分为两个截然不同的层次。
第一层叫做\grefdex{g:long-term-memory}{长期记忆}{长期记忆}
或\gref{g:persistent-memory}{持久记忆}，
我们把它用来存朋友的名字、
家里的地址、
以及八岁生日派对上小丑做了什么把我们吓得那么惨之类的事。
它的容量基本上是无限的，
但读取很慢——慢到帮不上我们应付饥饿的狮子和闹脾气的家人。

所以进化给了我们第二套系统，
叫做\grefdex{g:short-term-memory}{短期记忆}{短期记忆}
或\gref{g:working-memory}{工作记忆}。
它快得多，
但也小得多：~\cite{Mill1956}估计，成年人的工作记忆平均一次只能容纳 7±2 个条目。
这就是为什么\hreffoot{https://www.quora.com/Why-did-Bell-Labs-create-phone-numbers-of-7-digits-10-digits-Is-there-a-reason-that-dashes-and-brackets-are-used}{电话号码}
是 7 位或 8 位：
从前电话用的是拨号盘而不是键盘时，
这已经是大多数成年人能在拨号盘转上好几圈的时间里准确记住的最长数字串。

\begin{aside}{参与}
  工作记忆的大小有时被用来说明，
  为什么运动队往往有六个左右的成员，
  或者会被分成橄榄球里的前锋和后卫这样的小组。
  它也被用来说明，为什么会议只有到一定人数以内才有效率：
  如果二十个人试着讨论一件事，
  要么是三个会议同时在开，
  要么就是六个人在说话，其他所有人都在听。
  这个论点的意思是，人追踪同伴的能力受工作记忆大小的限制，
  但据我所知，
  这一联系从未被证实过。
\end{aside}

7±2 是教学里最重要的一个数字。
老师没法把信息直接放进学习者的长期记忆。
相反，
老师呈现的任何东西都先存进学习者的短期记忆，
只有在那里保持并反复演练（\secref{s:individual-strategies}）之后，才会转移到长期记忆。
如果老师呈现信息太多太快，
新信息就会在后一批信息被转移之前把前一批挤掉。

这是设计课程时使用概念图的方法之一：
它有助于确保学习者的短期记忆不会超载。
图画好之后，
老师挑选一个能放进短期记忆、并导向一次形成性评价的子部分（\figref{f:memory-photosynthesis}），\index{形成性评价}
然后为下一段课程再加一个子部分，依此类推。

\figpdf{figures/photosynthesis.pdf}{在课程设计中使用概念图}{f:memory-photosynthesis}{}

\begin{aside}{一起画概念图}
  下次开团队会议时，
  给每人发一张纸，
  让他们花几分钟为你们正在做的项目画一张自己的概念图。
  数到三，
  让所有人把自己的概念图展示给小组。
  随后的讨论，或许能帮大家明白
  为什么他们之前总是互相磕磕绊绊。
\end{aside}

请注意，这里呈现的简单记忆模型，在很大程度上已经被一个更精密的模型取代了。
在那个新模型里，短期记忆被拆成若干个模态存储
（例如视觉记忆与语言记忆），
每个存储都会做一些不自觉的预处理~\cite{Mill2016a}。
因此我们的呈现，是心智模型有助于学习和日常工作的一例。

\subsection*{模式识别}

近期研究表明，短期记忆的实际容量
可能低至 4±1 个条目~\cite{Dida2016}。
为了处理更大的信息量，
我们的头脑会形成\grefdex{g:chunking}{组块}{组块}。
例如，
我们大多数人记单词时，是把整个词当作一个条目，而不是一串字母。
同样，
纸牌或骰子上五个点的图案，是被当作一个整体记住的，
而不是五条分开的信息。

专家拥有比非专家更多、更大的组块，
也就是专家「看到」的是更大的模式，且有更多的模式可以用来匹配事物。
这使他们能在更高的层面推理，
并更快、更准地搜索信息。
然而，
如果我们认错了东西，组块也会误导我们：
新来的人有时确实能看到专家看过却漏掉的东西。

考虑到组块对思维如此重要，
人们很自然会想找出\hreffoot{https://en.wikipedia.org/wiki/Software\_design\_pattern}{设计模式}\index{设计模式}
并直接教它。
这些模式在许多领域（包括教学~\cite{Berg2012}）帮助熟手思考和相互交流，
但模式目录对新手来说太枯燥、太抽象，光靠他们自己没法理解。
话虽如此，
给少数模式起名字，似乎确实对教学有帮助，
主要是因为它给了学习者一套更丰富的词汇来思考和交流~\cite{Kuit2004,Byck2005,Saja2006}。
我们会在\secref{s:pck-programming}回到这一点。

\seclbl{成为专家}{s:memory-becoming-expert}

那么，一个人怎么成为专家？
「一万小时练习就能成」这个说法被广为引用，
但\hreffoot{http://www.goodlifeproject.com/podcast/anders-ericsson/}{大概并不是真的}：
把同一件事一遍遍重复，更可能巩固坏习惯，而不是提高表现。
真正管用的是做相似但又略有不同的事，
留心上什么管用、什么不管用，
然后根据这些反馈改变行为，从而越来越好。
这叫做\grefdex{g:deliberate-practice}{刻意练习}{刻意练习}
或\gref{g:reflective-practice}{反思性实践}，
一个常见的进程是人会经历三个阶段：

\begin{description}

\item[根据他人的反馈行动。]
  学习者可能写一篇关于自己暑假做了什么作文，
  然后从老师那里得到反馈，告诉他怎么改进。

\item[对他人的作品给出反馈。]
  学习者可能评论《哈利·波特》小说里的人物塑造，
  然后从老师那里得到关于他这篇评论的反馈。

\item[给自己反馈。]
  到了某个时候，
  学习者开始在做的时候，用自己已经积累起来的技能评论自己的作品。
  这样做比等别人的反馈快得多，
  于是熟练度会突然开始突飞猛进。

\end{description}

\begin{aside}{什么算刻意练习？}
  \cite{Macn2014}发现，
  「{\ldots}刻意练习解释了游戏表现方差的 26\%、
  音乐的 21\%、
  体育的 18\%、
  教育的 4\%，
  以及职业的不到 1\%。」
  不过，~\cite{Eric2016}对这一发现提出了批评，说：
  「把一个人生涯中任何类型的练习时间统统加起来，
  等于是在假设所有类型的练习活动对表现的影响都一样——这个假设
  {\ldots}与证据不符。」
  要有效，刻意练习需要一个清晰的表现目标，
  以及即时的、有信息量的反馈，
  而这两样本来就是老师应该努力去做到的。
\end{aside}

\seclbl{练习}{s:memory-exercises}

\exercise{概念图绘制}{两人一组}{30}

为你能在五分钟内讲清楚的东西画一张概念图。
和同伴交换，互相评论对方的图。
它们呈现的是概念，还是表面的细节？
同伴图里的哪些关系，你认为其实是概念？反过来呢？

\exercise{再画一次概念图}{小组}{20}

以 3—4 人为一组，
每人独立画一张概念图，展示自己对课堂上会发生什么的心智模型。
大家都画完后，
比较这些概念图。
你们的心智模型在哪里一致、哪里不一致？

\exercise{增强短期记忆}{个人}{5 分钟}

\cite{Cher2007}提出，
人们讨论事情时画图的主要原因，
是为自己的短期记忆扩容：
指着几分钟前画的一个歪歪扭扭的气泡，就能触发对好几分钟辩论的回忆。
在上一个练习里交换概念图时，
别人有多容易看懂你的图是什么意思？
如果你把图搁上一两天再看，你自己又会有多容易看懂？

\exercise{这有点自我指涉了，不是吗？}{全班}{30}

各自独立，
为「概念图」画一张概念图。
把你的概念图与别人画的比较。
大多数人放了哪些内容？
又有哪些显著的不同？

\exercise{注意你的盲区}{小组}{10}

Elizabeth Wickes 列出了
\hreffoot{https://twitter.com/elliewix/status/981285432922202113}{读懂下面这一行 Python 所需理解的所有东西}：

\begin{minted}{text}
answers = ['tuatara', 'tuataras', 'bus', "lick"]
\end{minted}

\begin{itemize}

\item
  内容两侧的方括号意味着我们在处理一个列表
  （而不是紧跟在某个东西右边的方括号，那是数据提取的记法）。

\item
  各元素之间用逗号分隔，逗号在引号外面、元素之间
  （而不是像引语里那样放在引号里面）。

\item
  每个元素都是一个字符串，
  我们从引号知道这一点。
  如果想要，这里也可以放数字或其他数据类型；
  之所以需要引号，是因为我们在处理字符串。

\item
  我们单引号和双引号混着用；
  只要每一对在单个字符串两侧是配对的，Python 并不在意。

\item
  每个逗号后面都跟了一个空格，
  这不是 Python 要求的，
  但我们为了可读性更愿意这样写。

\end{itemize}

这些细节中的每一个，专家都可能忽略。
以 3—4 人为一组，
从你最近教过或学过的一课里挑出一段同样短的内容，
把它拆解到这个细致程度。

\exercise{接下来教什么}{个人}{5}

回头看看\figref{f:memory-photosynthesis}里那张光合作用的概念图。
选中的区块里有多少概念和连线？
你会在课程的下一块里放什么？为什么？

\exercise{组块的力量}{个人}{5}

看\figref{f:memory-unchunked} 10 秒钟，
然后移开视线，试着用这些符号写出你的电话号码\footnote{
  感谢 Warren Code\index{Code, Warren}向我介绍这个例子。
}。
（用空格表示「0」。）
写完后，
看看\appref{s:chunking}里的另一种表示。
当模式被明确呈现出来之后，这些符号好记了多少？

\figpdfhere{figures/chunking-unchunked.pdf}{未组块的表示}{f:memory-unchunked}{}
